\section{Тестирование приложения}

\subsection{Подход к тестированию}

Функциональное тестирование разработанного приложения проведено в соответствии с рекомендациями методички (раздел~4.2) и разделено на критическое и углубленное.
Критическое тестирование охватывает основные сценарии работы пользователя при корректных входных данных; углубленное~--~поведение системы при некорректных, граничных и пограничных значениях.

В рамках проекта применены два метода функционального тестирования: метод <<чёрного ящика>> для проверки эндпоинтов REST API через инструмент Swagger UI и метод <<белого ящика>> для проверки логики алгоритма пересчёта результатов отбора через юнит-тесты класса \code{FullRecalculationCore}.

\subsection{Перечень аппаратных средств}

Аппаратные средства, использованные при тестировании, приведены в таблице~\ref{tab:hardware}.

\begin{longtblr}[
  caption = {Перечень аппаратных средств},
  label = {tab:hardware},
]{
  colspec = {|Q[c]|Q[l]|X|X|},
  width = \textwidth,
  hlines, vlines,
  rowhead = 1,
  row{1} = {halign=c},
}
№ & Роль & Аппаратная конфигурация & Программная конфигурация \\
1 & Сервер приложений & Apple M1, 16~ГБ ОЗУ & macOS 25, .NET 8 SDK, Docker \\
2 & Сервер базы данных & Тот же компьютер & PostgreSQL 18-alpine в контейнере Docker \\
3 & Рабочая станция & Apple M1, 16~ГБ ОЗУ & macOS 25, браузер Chrome 138 \\
\end{longtblr}

\subsection{Результаты функционального тестирования API}

Тестовые случаи для проверки эндпоинтов REST API и результаты их прохождения приведены в таблице~\ref{tab:apitests}.

\begin{longtblr}[
  caption = {Тестовые случаи для проверки REST API},
  label = {tab:apitests},
]{
  colspec = {|Q[c]|X|X|X|Q[c]|},
  width = \textwidth,
  hlines, vlines,
  rowhead = 1,
  row{1} = {halign=c},
}
№ & Название модуля & Описание тестового случая & Ожидаемый результат & Пройден? \\
1 & Аутентификация & POST \code{/auth/login} с корректными учётными данными & Возвращён JWT-токен, статус~200 & Да \\
2 & Аутентификация & POST \code{/auth/login} с неверным паролем & Статус~401, в \code{auth\_log} запись с неуспехом & Да \\
3 & Факультеты & POST \code{/faculties} от пользователя с ролью \code{SuperAdmin} & Факультет создан, статус~200 & Да \\
4 & Факультеты & POST \code{/faculties} от пользователя с ролью \code{DataViewer} & Статус~403 (Forbidden) & Да \\
5 & Факультеты & GET \code{/faculties/paged?page=1\&pageSize=10} & Возвращены первые 10 факультетов & Да \\
6 & Абитуриенты & POST \code{/applicants} с корректными данными & Абитуриент создан, статус~200 & Да \\
7 & Абитуриенты & DELETE \code{/applicants/\{id\}} & Создан запрос на удаление со статусом \code{pending}, абитуриент исключён из списков & Да \\
8 & Оценки & POST \code{/applicant\_evaluation\_values} со значением вне диапазона критерия & Статус~400, сообщение об ошибке & Да \\
9 & Оценки & POST дублирующего значения по той же паре (абитуриент, критерий) & Статус~409 (Conflict) благодаря уникальному ограничению & Да \\
10 & Отбор & POST \code{/selection/recalculate} & Таблица \code{selectedapplicants} перестроена & Да \\
11 & Отбор & GET \code{/selection/competition-lists/\{id\}/excel} & Возвращён файл \code{xlsx} с заголовком \code{Content-Disposition} & Да \\
12 & Аудит & POST \code{/audit/pending-deletions/\{id\}/confirm} & Запрос помечен \code{confirmed}, абитуриент скрыт окончательно & Да \\
13 & Аудит & POST \code{/audit/pending-deletions/\{id\}/reject} & Абитуриент возвращён, \code{validated=false}, выполнен пересчёт & Да \\
\end{longtblr}

\subsection{Тестирование алгоритма отбора}

Для проверки корректности работы алгоритма пересчёта результатов отбора разработан набор юнит-тестов, проверяющих следующие сценарии: распределение абитуриентов в одну категорию без вытеснения; вытеснение слабого кандидата при превышении квоты категории; вытеснение слабого кандидата при превышении общего плана конкурсного списка; корректную обработку нескольких заявок одного абитуриента в порядке приоритета предпочтения; работу с критериями типа \code{lower\_is\_better}; обработку категорий с пустой группой критериев.

Все разработанные тестовые случаи пройдены успешно, что подтверждает корректность реализации алгоритма.

\subsection{Перечень граничных и эквивалентных значений}

В соответствии с рекомендациями методички (раздел~4.2.2.6) для полей с числовыми значениями определены граничные и эквивалентные классы значений.
Результаты приведены в таблице~\ref{tab:boundary}.

\begin{longtblr}[
  caption = {Перечень граничных и эквивалентных значений},
  label = {tab:boundary},
]{
  colspec = {|X|X|X|X|},
  width = \textwidth,
  hlines, vlines,
  rowhead = 1,
  row{1} = {halign=c},
}
Название поля & Формат данных & Граничные значения & Эквивалентные значения \\
Значение оценки & \code{decimal} с двумя знаками после запятой & \code{minvalue}, \code{maxvalue} & \code{minvalue - 0.01}, среднее значение диапазона, \code{maxvalue + 0.01} \\
Квота категории & \code{integer} & 0, 1 & --10, 100 \\
План конкурсного списка & \code{integer} & 0, 1 & 100 \\
Приоритет защиты & \code{integer} & 1 & 2, 99 \\
Приоритет предпочтения & \code{integer} & 1 & 2, 100 \\
\end{longtblr}

\subsection{Выводы по тестированию}

По результатам функционального тестирования установлено, что разработанное приложение корректно реализует все сформулированные функциональные требования.
Все запланированные тестовые случаи (13 шт.) пройдены успешно.
Обнаруженные в процессе разработки ошибки были устранены до этапа итогового тестирования.

Алгоритм пересчёта результатов отбора, реализованный в классе \code{FullRecalculationCore}, продемонстрировал корректность работы на всех проверенных сценариях, включая граничные случаи с превышением квот и плана конкурсного списка.
